% Created 2023-02-26 Κυρ 18:24
% Intended LaTeX compiler: pdflatex
\documentclass[11pt]{article}
\usepackage[utf8]{inputenc}
\usepackage[T1]{fontenc}
\usepackage{graphicx}
\usepackage{longtable}
\usepackage{wrapfig}
\usepackage{rotating}
\usepackage[normalem]{ulem}
\usepackage{amsmath}
\usepackage{amssymb}
\usepackage{capt-of}
\usepackage{hyperref}
\usepackage{booktabs}
\usepackage{import}
\usepackage[LGR, T1]{fontenc}
\usepackage[greek, english]{babel}
\usepackage{alphabeta}
\usepackage{esint}
\usepackage{mathtools}
\usepackage{esdiff}
\usepackage{makeidx}
\usepackage{glossaries}
\usepackage{newfloat}
\usepackage{minted}
\usepackage{chemfig}
\usepackage{svg}
\usepackage[a4paper, margin=3cm]{geometry}
\author{Βιδιάνος Γιαννίτσης}
\date{\today}
\title{Παραγωγή Ενέργειας στην μονάδα και Ενεργειακή Ολοκλήρωση}
\hypersetup{
 pdfauthor={Βιδιάνος Γιαννίτσης},
 pdftitle={Παραγωγή Ενέργειας στην μονάδα και Ενεργειακή Ολοκλήρωση},
 pdfkeywords={},
 pdfsubject={},
 pdfcreator={Emacs 28.2 (Org mode 9.5.5)}, 
 pdflang={English}}
\makeatletter
\newcommand{\citeprocitem}[2]{\hyper@linkstart{cite}{citeproc_bib_item_#1}#2\hyper@linkend}
\makeatother

\usepackage[notquote]{hanging}
\begin{document}

\maketitle
\tableofcontents

\renewcommand{\abstractname}{Περίληψη}
\renewcommand{\tablename}{Πίνακας}
\renewcommand{\figurename}{Σχήμα}
\renewcommand\listingscaption{Κώδικας}

\section{Εισαγωγή}
\label{sec:org642b3db}
Στο αρχείο αυτό θα γίνει μία συνοπτική παρουσίαση της πρόοδου της εργασίας μας από την τελική εξέταση του προηγούμενου εξαμήνου \href{https://github.com/Vidianos-Giannitsis/Process-Design/blob/master/Final\_exam\_files/olive\_kernel\_glycerol\_cyclopentanone\_omada\_27\_final.pdf}{δείτε εδώ} μέχρι τώρα. Ο στόχος μας ήταν να μελετήσουμε περισσότερο τα ενεργειακά της διεργασίας στα οποία δεν είχε δωθεί ιδιαίτερη έμφαση στο προηγούμενο εξάμηνο. Η δουλειά που έγινε μπορεί να χωριστεί σε 2 βασικά κομμάτια. Την ανάπτυξη του \href{https://github.com/Vidianos-Giannitsis/Process-Design/blob/master/Final\_exam\_files/block\_300\_info.pdf}{block 300} της διεργασίας, το οποίο αποτελεί την μονάδα εκμετάλλευσης της λιγνίνης που ανακτάται στο block 100 και το οποίο προχώρησε από έναν απλό καυστήρα σε ένα ολοκληρωμένο κύκλο Rankine με μονάδες αναθέρμανσης, αναγέννησης και τηλεθέρμανσης καθώς και την ενεργειακή ολοκλήρωση της διεργασίας και την δημιουργία του μεγάλου σύνθετου γραφήματος της.

\section{Ανάπτυξη του block 300 της διεργασίας}
\label{sec:org845a65c}
To block 300 ξεκίνησε στο προηγούμενο εξάμηνο ως ένας απλός καυστήρας από τον οποίο παραγόταν ατμός υψηλής πίεσης ως θερμαντικό μέσο της εγκατάστασης. Καθώς η μονάδα επεξεργάζεται πολύ μεγάλες ποσότητες βιομάζας, η χρήση όλης αυτής της λιγνίνης ως θερμαντικό μέσο δεν είναι ιδιαίτερα ουσιώδης καθώς θα παραχθεί πολύ μεγαλύτερη ποσότητα ατμού από αυτήν που χρειάζεται η εγκατάσταση (συγκεκριμένα, η δυναμικότητα της είναι 81 MW θερμότητας). Για αυτό, μελετήθηκε η ιδέα της δημιουργίας ενός κύκλου Rankine το οποίο θα χρησιμοποιεί την θερμική ενέργεια των καυσαερίων της λιγνίνης για ηλεκτροπαραγωγή. Για να έχει υψηλή απόδοση το κύκλο, υπάρχουν μονάδες αναθέρμανσης και αναγέννησης και χρησιμοποιούνται 4 διαφορετικές βαθμίδες ατμού. Η πρώτη είναι ο ατμός υψηλής πίεσης του λέβητα (10 MPa), η δεύτερη είναι επίσης ατμός σχετικά υψηλής πίεσης (4 MPa) ο οποίος μπορεί να χρησιμοποιηθεί ως θερμανικό μέσο για την υπόλοιπη εγκατάσταση. Η τρίτη βαθμίδα είναι ατμός μέσης πίεσης (2.5 MPa) ο οποίος χρησιμοποιείται για την μονάδα αναγέννησης για να βελτιώσει την απόδοση της εγκατάστασης ενώ η τελευταία βαθμίδα είναι το νερό χαμηλής πίεσης για να κλείσει το κύκλο, το οποίο είναι στα 0.1 bar.

Για το κύκλο, έγινε πρώτα ένας σχεδιασμός ως "proof of concept" στο χαρτί όπου έγιναν οι βασικοί υπολογισμοί των ρευμάτων του κύκλου και εξετάστηκε η σκοπιμότητα ενός τέτοιου κύκλου και κατά πόσο μπορεί να καλύψει όλες τις ενεργειακές ανάγκες της εγκατάστασης. Αυτό βασίστηκε στο παρακάτω διάγραμμα ροής
\begin{figure}[htbp]
\centering
\includegraphics[width=.9\linewidth]{/home/vidianos/Documents/7o_εξάμηνο/Σχεδιασμός_Ι/Project/git_repo/5th_exam_files/energy_production_and_integration_5th_exam-org-img/rankine_cycle.png}
\caption{Πρόχειρο διάγραμμα ροής της διεργασίας}
\end{figure}
Το πρόβλημα λύθηκε σε αυτήν την φάση θεωρώντας ιδανική λειτουργία των στρεπτών, ότι η διαθέσιμη θερμική ενέργεια είναι αυτή των καυσαερίων της λιγνίνης και η απαιτούμενη θερμική ενέργεια αυτή όλων των υπόλοιπων σημείων της εγκατάστασης χωρίς ενεργειακή ολοκλήρωση. Υπολογίστηκε πως η διεργασία είναι εφικτή, για αυτό ως επόμενο στάδιο, αυτή προσομοιώθηκε στο Aspen, λαμβάνοντας υπόψην τις μη ιδανικότητες και βελτιστοποιώντας τις συνθήκες για τις ενεργειακές απαιτήσεις με την ενεργειακή ολοκλήρωση.

Το αντίστοιχο διάγραμμα ροής δημιουργημένο στο Aspen φαίνεται παρακάτω
\begin{center}
\includegraphics[width=.9\linewidth]{Ανάπτυξη_του_block_300_της_διεργασίας/2023-02-26_18-23-52_screenshot.png}
\end{center}

Λεπτομέρειες της προσομοίωσης στο Aspen φαίνονται στο \href{https://github.com/Vidianos-Giannitsis/Process-Design/blob/master/Aspen/lignin\_rankine\_cycle.org}{σχετικό αρχείο} αλλά αξίζει να αναφερθεί πως η δυναμικότητα του είναι 24.2 MW ηλεκτρισμού και 5.2 MW θερμότητας η οποία διατίθεται ως υπέρθερμος ατμός στους 364.8 \(^oC\) και 40 bar. Όπως θα αναφερθεί και παρακάτω, το θερμικό φορτίο αυτό είναι επαρκές για την θέρμανση όλων των ρευμάτων της διεργασίας.

\section{Ενεργειακή Ολοκλήρωση - Μεγάλο Σύνθετο Γράφημα}
\label{sec:org9eba468}
Για την ενεργειακή ολοκλήρωση της διεργασίας, πρέπει αρχικά να αναγνωριστούν όλα τα απαιτούμενα θερμά και ψυχρά ρεύματα της διεργασίας. Αυτά φαίνονται στον παρακάτω πίνακα

\begin{table}[htbp]
\caption{Χαρακτηρισμός των Ρευμάτων}
\centering
\begin{tabular}{llrrrrr}
Ρεύμα & Χαρακτηρισμός & Τ\textsubscript{αρχ} & Τ\textsubscript{τελ} & Cp J/mol K & CP MJ/K h & ΔΗ MJ/h\\
\hline
FeedSteam & Ψυχρό & 20 & 232 & 75.6 & 47.872 & 3.58e+4\\
CellLig & Θερμό & 232 & 50 & 230.41 & 13.36 & 2430.95\\
Bleach & Ψυχρό & 20 & 70 & 75.6 & 4.2049 & 210.25\\
SachCell & Ψυχρό & 50 & 51.4 & 80.78 & 0.06 & 0.0842\\
Gluc & Θερμό & 50 & 30 & 79.3 & 53.08 & 1061.6\\
RProd & Ψυχρό & 30 & 150 & 75.4 & 58.38 & 3.53e+4\\
Xylose & Ψυχρό & 232 & 243 & 180 & 4.75 & 52.25\\
Furf & Θερμό & 242 & 160 & 96.1 & 10.14 & 831.48\\
\end{tabular}
\end{table}

Στον πίνακα αυτόν δεν αναφέρονται τα ρεύματα των τελικών προιόντων τα οποία θα μπορούσαν να παίξουν τον ρόλο θερμών ρευμάτων και να ψυχθούν μέχρι θερμοκρασία περιβάλλοντος για την αποθήκευση τους για να δούμε πρώτα τις απαιτήσεις που υπάρχουν χωρίς αυτά. Σε επόμενο στάδιο, μπορούν να συμπεριληφθούν και αυτά.

Με κάποιους απλούς υπολογισμούς, που παρουσιάζονται \href{https://github.com/Vidianos-Giannitsis/Process-Design/blob/master/energy\_integration.org}{εδώ}, προέκυψε ο σχετικός ενεργειακός καταρράκτης και το μεγάλο σύνθετο γράφημα της διεργασίας που φαίνεται παρακάτω.

\begin{figure}[htbp]
\centering
\includegraphics[width=.9\linewidth]{Ενεργειακή_Ολοκλήρωση_-_Μεγάλο_Σύνθετο_Γράφημα/2023-02-26_18-23-13_screenshot.png}
\caption{Μεγάλο σύνθετο γράφημα της διεργασίας}
\end{figure}

Από εδώ φαίνονται οι ενεργειακές απαιτήσεις της διεργασίας ως

Απαίτηση σε θερμή παροχή: 13613.4 MJ/h ή 3.78 MW και απαίτηση σε ψυχρή παροχή: 10 MJ/h ή 2.78 kW. Όπως φαίνεται και από τα νούμερα αυτά, η διαθέσιμη θερμή παροχή που παράγεται στην μονάδα από το block 300 επαρκεί για να καλύψει όλες τις υπόλοιπες ενεργειακές απαιτήσεις της διεργασίας. Η ψυχρή παροχή που απαιτείται εδώ είναι 2.78 kW, βέβαια αξίζει να αναφερθεί ότι υπάρχουν και άλλες απαιτήσεις σε διεργασίες που δεν έχουν ληφθεί υπόψην στο γράφημα αυτό (πχ αντιδραστήρες).

Αξίζει να συζητηθούν και ορισμένα \href{https://github.com/Vidianos-Giannitsis/Process-Design/blob/master/energy\_integration.org\#\%CF\%83\%CF\%87\%CF\%8C\%CE\%BB\%CE\%B9\%CE\%B1-\%CE\%B3\%CE\%B9\%CE\%B1-\%CF\%84\%CE\%B7\%CE\%BD-\%CE\%BF\%CE\%BB\%CE\%BF\%CE\%BA\%CE\%BB\%CE\%AE\%CF\%81\%CF\%89\%CF\%83\%CE\%B7-\%CE\%B4\%CE\%B9\%CE\%AC\%CF\%86\%CE\%BF\%CF\%81\%CF\%89\%CE\%BD-\%CE\%BA\%CE\%BF\%CE\%BC\%CE\%BC\%CE\%B1\%CF\%84\%CE\%B9\%CF\%8E\%CE\%BD}{σχόλια} που έχουν γίνει πάνω στην ολοκλήρωση της υπόλοιπης διεργασίας με τα ρεύματα αυτά (αντιδραστήρες-αποστακτικές).

\section{Απορίες}
\label{sec:org845481b}
\begin{enumerate}
\item Πως μπορούν να προσεγγιστούν οι απαιτήσεις της εγκατάστασης σε ηλεκτρική ενέργεια; Στο κύκλο που σχεδιάστηκε, οι βαθμίδες πίεσεις επιλέχθηκαν αυθαίρετα για να είναι βολικά τα σημεία από άποψης θερμότητας και η ηλεκτρική ενέργεια προέκυψε από αυτά. Αλλά υπήρχε η σκέψη να υπολογιστούν κάπως αυτές οι απαιτήσεις αν γίνεται.

\item Πως ακριβώς μεταφράζεται το μεγάλο σύνθετο γράφημα στο Aspen? Αυτό που σκεφτόμουν είναι ότι στην ουσία βάζουμε τα ρεύματα σε κάθε βαθμίδα του ενεργειακού καταρράκτη να εναλλάξουν μεταξύ τους θερμότητα και συμπληρώνουμε σε κάθε περίπτωση την θερμή παροχή που απαιτείται αλλά ήθελα να ρωτήσω επειδή δεν είμαι σίγουρος.

\item Πως μπορεί να γίνει (αν μπορεί) η ολοκλήρωση αντιδραστήρων και αποστακτικών στηλών στο Aspen? Είχαμε δει στο μάθημα πως με βάση το μεγάλο σύνθετο γράφημα, μπορεί να γίνει ολοκλήρωση αυτών, αλλά ήθελα να ρωτήσω αν υπάρχει τρόπος αυτό να περαστεί στο Aspen μιας και τουλάχιστον όπως το έχουμε χρησιμοποιήσει, δεν έχουμε δει κάπως τρόπο να ορίσουμε το ρεύμα εναλλαγής των διατάξεων αυτών.
\end{enumerate}
\end{document}